\documentclass[12pt]{article}
\usepackage[italian]{babel}
\usepackage[T1]{fontenc}
\usepackage[utf8]{inputenc}
\usepackage{amsmath}
\usepackage{amssymb}
\usepackage{bm}
\usepackage{amsthm}
\usepackage{xspace}
\usepackage{geometry}
\geometry{a4paper, top=1.9cm, bottom=1.9cm, left=2.2cm, right=2.2cm}

\usepackage[colorlinks,allcolors=black]{hyperref}

\linespread{1.2}

\newcommand{\N}{\mathbb{N}}
\newcommand{\R}{\mathbb{R}}
\newcommand{\Z}{\mathbb{Z}}
\newcommand{\E}{\mathbb{E}}
\newcommand{\val}{\operatorname{val}}
\newcommand{\CSP}{\mathtt{CSP}\xspace}
\renewcommand{\P}{\mathbf{P}\xspace}
\newcommand{\NP}{\mathbf{NP}\xspace}
\newcommand{\PCP}{\mathbf{PCP}\xspace}
\newcommand{\F}{\mathbb{F}}

\begin{document}
\pagestyle{empty}

\begin{center}
  {\small\scshape Università di Pisa -- Dipartimento di Matematica \par}
  {\small Corso di Laurea Triennale in Matematica \par}
  \vspace{0.8em}
  {\Large\bfseries Teorema PCP e inapprossimabilità \par}
\end{center}

\vspace{0.8em}

\noindent
\begin{minipage}[t]{0.47\textwidth}
  {\large Relatori: \par}
  \vspace{0.2em}
  {\bfseries\large Prof. Alessandro Panconesi \par}
  {\bfseries\large Prof.ssa Alessandra Caraceni \par}
\end{minipage}%
\hfill
\begin{minipage}[t]{0.47\textwidth}\raggedleft
  {\large Candidato: \par}
  \vspace{0.2em}
  {\bfseries\large Lorenzo Ferrari \par}
\end{minipage}

\vspace{0.8em}
\hrule height 0.4pt
\vspace{0.8em}

\begin{center}
  {\large\bfseries Abstract \par}
\end{center}
\vspace{0.2em}

Lo scopo di questa tesi è lo studio del Teorema PCP (\emph{Probabilistically
Checkable Proofs}), risultato fondamentale in informatica teorica, dimostrato
in una prima versione nel 1992 da Arora e Safra e completato lo stesso anno da
Arora, Lund, Motwani, Sudan e Szegedy.

Una dimostrazione si legge tipicamente dall'inizio alla fine per assicurarsi
che tutti i passaggi logici siano validi. Un singolo errore, infatti, potrebbe
invalidare l'intero argomento.
Se tuttavia ci venisse fornita una dimostrazione molto lunga di un teorema
complicato, ci piacerebbe adottare la seguente procedura di verifica: ``lette
cento righe scelte a caso, accettiamo la dimostrazione come valida se e soltanto
se questa non contiene errori''.
Il Teorema PCP afferma che esiste un modo di scrivere dimostrazioni per cui la
verifica descritta ha successo con alta probabilità:
dimostrazioni corrette vengono certamente accettate e, al tempo stesso,
ogni dimostrazione fallace viene rifiutata con probabilità (diciamo) maggiore
del 99\%.
Questa tesi intende presentare in maniera precisa il Teorema PCP, analizzarne le
conseguenze per l'inapprossimabilità di problemi di ottimizzazione combinatoria
e infine percorrerne la dimostrazione combinatoria di Irit Dinur.

Nel primo capitolo introduciamo le nozioni base della teoria della
complessità a partire dalle classi $\P$ e $\NP$, che catturano
rispettivamente i problemi che si risolvono efficientemente e i problemi le cui
soluzioni si verificano efficientemente.
Descriviamo i problemi $\NP$-completi $\mathtt{SAT}$, $\mathtt{3SAT}$,
$\mathtt{3COL}$ e soprattutto $\mathtt{QUADEQ}$: il problema di determinare se
un sistema di equazioni quadratiche su $\F_2$ ha soluzione.
Formalizziamo poi il modello di verifica probabilistica di dimostrazioni ed
enunciamo il Teorema PCP.
Come primo risultato significativo, presentiamo la dimostrazione di una versione
debole del Teorema PCP, nota come Teorema PCP esponenziale.
La dimostrazione fa uso del codice di Walsh-Hadamard per rappresentare
dimostrazioni come funzioni lineari. Sarà in particolare necessario un test di
linearità per funzioni $\F_2^n \to \F_2$, la cui probabilità di successo si
stima mediante analisi di Fourier discreta sull'ipercubo
$(\Z / 2\Z)^n$.

Nel secondo capitolo esploriamo il legame tra verifica probabilistica
di dimostrazioni e inapprossimabilità di problemi di ottimizzazione
combinatoria. Generalizziamo $\mathtt{3SAT}$ e $\mathtt{3COL}$ introducendo
i \emph{Constraint Satisfaction Problems} ($\mathtt{CSP}$) e mostriamo che il
Teorema PCP ammette una formulazione equivalente intrinsecamente legata
all'inapprossimabilità di $\CSP$.
Come conseguenza diretta, non esiste un algoritmo efficiente in grado di
approssimare il valore di soddisfacibilità di istanze di $\CSP$ oltre una
certa soglia costante a meno che $\P = \NP$.
Concludiamo con due risultati di inapprossimabilità storicamente rilevanti.
Il primo, per il problema della cricca massima, afferma che non si può trovare
efficientemente neanche una cricca grande un miliardesimo della cricca massima,
a meno che $\P = \NP$.
Il secondo, diretta conseguenza del Teorema PCP, afferma che sotto la stessa
ipotesi non si può approssimare efficientemente \texttt{MAX-3SAT} oltre una
certa soglia costante.
Grazie ai lavori precedenti, ciò sbloccò in un colpo solo risultati di
inapprossimabilità per molti altri problemi.

Il terzo capitolo è dedicato ai grafi expander e alle loro proprietà
pseudo-casuali. Per esempio, per ogni nodo $v$ di un expander con $n$ nodi,
la distribuzione del $k$-esimo nodo di una passeggiata aleatoria che parte da
$v$ è molto vicina a quella uniforme per $k = O(\log n)$.
Mostriamo inoltre che, per molte applicazioni, i nodi di una passeggiata
aleatoria di lunghezza $k$ funzionano bene quasi quanto $k$ nodi scelti a caso.
Ciò permette il cosiddetto ``riciclo di bit casuali'', che incluso nella
dimostrazione di inapprossimabilità per il problema della cricca massima
fornisce un risultato ancora migliore.
Ci interessiamo infine all'esistenza di expander e alla loro costruzione.

Nell'ultimo capitolo percorriamo la dimostrazione combinatoria
del Teorema PCP proposta da Dinur nel 2006 basandoci sulla trattazione di
Arora e Barak.
La dimostrazione si basa sull'\emph{amplificazione del gap}, ottenuta grazie
all'impiego degli expander e alle loro proprietà pseudo-casuali.
Tra gli ingredienti della dimostrazione figura anche il Teorema PCP
esponenziale, dimostrato nel primo capitolo.

\end{document}
